\section{Gaussian Space --- Hermite --- Truncation}
%%% generated by the blueprint pipeline from TCSlib.BooleanAnalysis.GaussianSpace.Hermite.Truncation
\subsection*{Gaussian Space --- Hermite --- Truncation: Introduction}

Declarations of the Lean module \texttt{TCSlib.BooleanAnalysis.GaussianSpace.Hermite.Truncation}.

\subsection*{Gaussian Space --- Hermite --- Truncation: Results}

\begin{definition}[Hermite degree indices]
\lean{BooleanAnalysis.GaussianSpace.hermiteDegreeIndices}
\label{BooleanAnalysis.GaussianSpace.hermiteDegreeIndices}
\leanok
For a natural number $n$ and a natural number $k$, let $I_{n,k}$ be the finite set of
multi-indices $\alpha:\{0,\ldots,n-1\}\to\mathbb{N}$ whose total degree satisfies
$\sum_{i=0}^{n-1}\alpha_i\le k$.
\end{definition}

\begin{lemma}[Characterization of Hermite degree indices]
\lean{BooleanAnalysis.GaussianSpace.mem_hermiteDegreeIndices}
\label{BooleanAnalysis.GaussianSpace.mem_hermiteDegreeIndices}
\leanok
\uses{BooleanAnalysis.GaussianSpace.hermiteDegreeIndices}
\difficulty{4}
Let $n$ and $k$ be natural numbers. An $n$-tuple $\alpha=(\alpha_1,\ldots,\alpha_n)$ of
natural numbers belongs to the set of Hermite degree indices of order $k$ exactly when
$\sum_{i=1}^n\alpha_i\le k$.
\end{lemma}

\begin{lemma}[Exhaustion by Hermite degree indices]
\lean{BooleanAnalysis.GaussianSpace.hermiteDegreeIndices_tendsto}
\label{BooleanAnalysis.GaussianSpace.hermiteDegreeIndices_tendsto}
\leanok
\uses{BooleanAnalysis.GaussianSpace.hermiteDegreeIndices, BooleanAnalysis.GaussianSpace.mem_hermiteDegreeIndices}
\difficulty{4}
For each natural number $n$, let $I_k$ be the set of multi-indices
$\alpha:\{1,\ldots,n\}\to\mathbb N$ whose total degree is at most $k$. Every finite
collection of such multi-indices is contained in $I_k$ for all sufficiently large $k$.
\end{lemma}

\begin{lemma}[Sum over bounded Hermite indices]
\lean{BooleanAnalysis.GaussianSpace.sum_hermiteDegreeIndices}
\label{BooleanAnalysis.GaussianSpace.sum_hermiteDegreeIndices}
\leanok
\uses{BooleanAnalysis.GaussianSpace.hermiteDegreeIndices}
\difficulty{3}
Let $n,k\in\mathbb N$, let $E$ be an additive commutative monoid, and assign an element
$v(\alpha)\in E$ to each function $\alpha\colon\{1,\ldots,n\}\to\mathbb N$. The sum of
$v(\alpha)$ over all such indices with $\sum_i\alpha(i)\le k$ equals the sum over all
coordinate choices $a(i)\in\{0,\ldots,k\}$, retaining $v(a)$ when $\sum_i a(i)\le k$ and
contributing zero otherwise.
\end{lemma}

\begin{lemma}[Finite Hermite sum representative]
\lean{BooleanAnalysis.GaussianSpace.coeFn_sum_hermiteL2}
\label{BooleanAnalysis.GaussianSpace.coeFn_sum_hermiteL2}
\leanok
\uses{BooleanAnalysis.GaussianSpace.hermiteFunction, BooleanAnalysis.GaussianSpace.hermiteFunction_memLp_two, BooleanAnalysis.GaussianSpace.hermiteL2}
\difficulty{4}
Let $n$ be a natural number, let $S$ be a finite set of multi-indices in $\mathbb N^n$,
and assign a real coefficient $c_\alpha$ to each multi-index. Under standard Gaussian
measure on $\mathbb R^n$, the function represented in $L^2$ by the finite sum of the
corresponding Hermite functions, with coefficients $c_\alpha$, agrees almost everywhere
with the pointwise sum $x\mapsto\sum_{\alpha\in S}c_\alpha H_\alpha(x)$, where
$H_\alpha$ denotes the multivariate Hermite function.
\end{lemma}

\begin{theorem}[Convergence of Hermite degree truncations]
\lean{BooleanAnalysis.GaussianSpace.hermite_degree_truncation_tendsto}
\label{BooleanAnalysis.GaussianSpace.hermite_degree_truncation_tendsto}
\leanok
\uses{BooleanAnalysis.GaussianSpace.coeFn_sum_hermiteL2, BooleanAnalysis.GaussianSpace.hermiteCoeff, BooleanAnalysis.GaussianSpace.hermiteDegreeIndices, BooleanAnalysis.GaussianSpace.hermiteDegreeIndices_tendsto, BooleanAnalysis.GaussianSpace.hermiteFunction, BooleanAnalysis.GaussianSpace.hermiteL2, BooleanAnalysis.GaussianSpace.hermite_expansion, BooleanAnalysis.GaussianSpace.sum_hermiteDegreeIndices}
\difficulty{5}
Let $n$ be a natural number and let $f:\mathbb{R}^n\to\mathbb{R}$ belong to $L^2$ for
the product of $n$ standard Gaussian measures. For each natural number $k$, form the sum
over multi-indices whose coordinates are at most $k$, retaining only those of total
degree at most $k$, with each term the Hermite coefficient of $f$ times the
corresponding normalized multivariate Hermite function. The $L^2$ distance between $f$
and this sum tends to zero as $k\to\infty$.
\end{theorem}
